% capitulos/05-resultados.tex

Lorem ipsum dolor sit amet, consectetur adipiscing elit, sed do eiusmod
tempor incididunt ut labore et dolore magna aliqua. Os resultados obtidos
nas três campanhas de campo revelaram padrões consistentes de distribuição
de mercúrio total (HgT) nos compartimentos ambientais investigados
\cite{exemplo:artigo:2024}.

\section{Concentrações de Mercúrio no Solo}

As concentrações de HgT nos solos variaram entre 12,4 e 187,3 ng g$^{-1}$
(peso seco), com média de $64{,}2 \pm 31{,}7$ ng g$^{-1}$ para o
horizonte superficial (0--10 cm). A \autoref{tab:resultados-hg} apresenta
as concentrações médias por horizonte e por campanha.

\begin{table}[H]
  \centering
  \caption{Concentrações médias de HgT (ng g$^{-1}$, peso seco) nos
           horizontes de solo por campanha de coleta.
           Letras diferentes indicam diferença significativa entre
           campanhas (Tukey, $p < 0{,}05$).}
  \label{tab:resultados-hg}
  \begin{tabularx}{\textwidth}{Xccc}
    \toprule
    \textbf{Horizonte (cm)}
      & \textbf{Campanha 1 (seca)}
      & \textbf{Campanha 2 (transição)}
      & \textbf{Campanha 3 (chuvosa)} \\
    \midrule
    0--10  & $64{,}2 \pm 31{,}7^{\,a}$ & $71{,}5 \pm 28{,}4^{\,a}$ & $58{,}3 \pm 25{,}1^{\,b}$ \\
    \thinrule
    10--20 & $48{,}1 \pm 22{,}3^{\,a}$ & $52{,}7 \pm 19{,}8^{\,a}$ & $44{,}9 \pm 18{,}6^{\,a}$ \\
    \thinrule
    20--30 & $31{,}4 \pm 15{,}9^{\,a}$ & $34{,}2 \pm 14{,}1^{\,a}$ & $29{,}8 \pm 13{,}7^{\,a}$ \\
    \bottomrule
  \end{tabularx}
  \fonte{Dados da pesquisa (2025).}
\end{table}

\section{Relação entre HgT e Matéria Orgânica do Solo}

A regressão linear simples entre HgT e COT revelou correlação positiva
e significativa ($r^2 = 0{,}72$; $p < 0{,}001$), indicando que a matéria
orgânica do solo é um preditor relevante da concentração de mercúrio,
conforme relatado por \citeonline{exemplo:artigo:2024} e
\citeonline{exemplo:capitulo:2019}.

\begin{figure}[H]
  \centering
  \fbox{\rule{0pt}{7cm}\rule{12cm}{0pt}}
  \caption{Regressão linear entre concentração de HgT (ng g$^{-1}$)
           e carbono orgânico total -- COT (\%) nos solos da área de
           estudo (ESECAE, DF, $n = 270$). A linha sólida representa
           o ajuste linear e a área sombreada o intervalo de confiança
           de 95\%. Fonte: dados da pesquisa (2025).}
  \label{fig:regressao-hg-cot}
\end{figure}

\section{Discussão}

Os resultados demonstram que a distribuição vertical de HgT no perfil
do solo segue padrão de decréscimo com a profundidade
(\autoref{tab:resultados-hg}), corroborando achados de
\citeonline{exemplo:tese:2023}. Esse padrão é atribuído à maior
concentração de matéria orgânica nos horizontes superficiais, que
funcionam como sorvedouro geoquímico de HgT \cite{exemplo:livro:2020}.

A variação sazonal observada entre as campanhas 1 e 3
(\autoref{tab:resultados-hg}) pode ser explicada pela lixiviação
do mercúrio para horizontes mais profundos durante o período chuvoso,
fenômeno descrito por \citeonline{exemplo:artigo:2024} para solos
tropicais com alta pluviosidade. A forte correlação entre HgT e COT
(\autoref{fig:regressao-hg-cot}) reforça o papel da matéria orgânica
como principal controlador da retenção de mercúrio nesse tipo de solo
\cite{exemplo:capitulo:2019}.
